\documentclass[12pt]{article}
\usepackage[utf8]{inputenc}
\usepackage[T1]{fontenc}
\usepackage[margin=1in]{geometry}
\usepackage{setspace}
\usepackage{xcolor}
\usepackage[colorlinks=true, urlcolor=blue!60!black, linkcolor=blue!60!black]{hyperref}
\setlength{\parindent}{2em}
\setlength{\parskip}{1.25em}
\renewcommand{\baselinestretch}{1.0}
\begin{document}
\begin{center}
{\LARGE \textit{Humphrey's Executor v.\ United States}, 295 U.S.\ 602 (1935)\par}
\vspace{0.8em}
{\large Case Brief --- Four-Page Edition\par}
\end{center}
\vspace{1em}

\section{Memory Jogger}

A unanimous Court held that Congress could limit the President to ``for-cause'' removal of Federal Trade Commissioners, because the FTC performed ``quasi-legislative'' and ``quasi-judicial'' duties rather than ``purely executive'' ones. The decision was overruled in 2026 by \textit{Trump v.\ Slaughter} (\href{https://supreme.justia.com/cases/federal/us/295/602/}{Justia}).

\section{Detailed Case Facts}

President Hoover renominated William E.\ Humphrey, a conservative Republican, to the FTC in 1931 for a seven-year term ending in 1938. The FTC Act of 1914 created a five-member commission with staggered seven-year terms, required partisan balance, and provided that a commissioner ``may be removed by the President for inefficiency, neglect of duty, or malfeasance in office.'' The Commission prevented unfair methods of competition through complaints, hearings, and cease-and-desist orders, investigated and reported to Congress, and served as a master in chancery for courts. In 1933 President Franklin Roosevelt asked Humphrey to resign so the Administration's aims could be carried out ``with personnel of my own selection,'' later writing that ``your mind and my mind go along together'' on neither policy nor administration. When Humphrey refused, Roosevelt removed him on October 7, 1933, citing none of the statutory causes. Humphrey insisted he remained a commissioner until his death in February 1934.

\section{Procedural History}

Humphrey's executor, Samuel Rathbun, sued in the Court of Claims for his unpaid salary. That court certified two questions: whether the FTC Act limited removal to the listed causes, and if so, whether that limit was constitutional. The Supreme Court answered both ``Yes'' on May 27, 1935---the same day it struck down the National Industrial Recovery Act in \textit{Schechter Poultry}.

\section{Judicial Votes}

\textbf{9--0}, both certified questions answered in the affirmative.
\begin{itemize}
\item \textbf{Opinion of the Court:} Justice Sutherland, joined by Chief Justice Hughes and Justices Van Devanter, Brandeis, Butler, Stone, Roberts, and Cardozo.
\item \textbf{Separate statement:} Justice McReynolds agreed, referring to his \textit{Myers} dissent.
\item \textbf{Dissents:} None.
\end{itemize}

\section{Holding}

When Congress fixes a definite term for members of an independent, nonpartisan commission and enumerates causes for removal, it bars removal during the term except for those causes. Constitutionally, whether the President's removal power may be limited ``will depend upon the character of the office'': \textit{Myers v.\ United States} protects illimitable removal only of purely executive officers, while Congress may impose for-cause protection on officers of agencies performing quasi-legislative and quasi-judicial functions. The Court left the ``field of doubt'' between those categories for future cases.

\section{Analysis of Opinions}

\textbf{Majority (Sutherland): statute.} \textit{Shurtleff v.\ United States} did not control because that statute fixed no term, so reading its causes as exclusive would have created life tenure. The FTC Act's seven-year terms and its legislative history showed Congress wanted an expert body ``independent of executive authority, except in its selection.''

\textbf{Majority (Sutherland): Constitution.} \textit{Myers} involved a postmaster, a purely executive officer; its broader language was dicta. The FTC carried out legislative policies, reported to Congress, adjudicated, and aided courts; any executive function was incidental. Congress's power to protect such officers flows from the separation of powers itself: an officer serving at another's pleasure cannot act independently, and giving the President illimitable removal power over the aids of the other branches would let the Executive coerce them. Sutherland reread the Decision of 1789 and \textit{Marbury} as distinguishing purely executive officers from those with judicial duties.

\textbf{McReynolds.} He joined the answers and pointed to his \textit{Myers} dissent, which denied the President any inherent, illimitable removal power over statutory officers.

\section{Examples: Future Applications}

\begin{enumerate}
\item \textbf{Same-Side: Claims Commission.} A multi-member commission that mainly adjudicates statutory claims performs quasi-judicial duties, so Congress may bar at-will removal (as in \textit{Wiener v.\ United States}, 1958).
\item \textbf{Same-Side: Bipartisan Expert Board.} A five-member board with staggered terms, partisan balance, and the FTC's removal language fit the classic blueprint---though after \textit{Slaughter}, multi-member agencies that enforce mandates against private parties are subject to at-will removal.
\item \textbf{Opposite-Side: Cabinet Secretary.} A statute protecting the Secretary of State from removal except for malfeasance would restrict a purely executive officer and be unconstitutional.
\item \textbf{Opposite-Side: Single Director.} A lone director with rulemaking and enforcement power falls outside \textit{Humphrey's Executor} (\textit{Seila Law v.\ CFPB}, 2020).
\item \textbf{Fence-Sitter: Federal Reserve.} The Board of Governors' central-banking lineage placed it in contested middle ground; on the day it decided \textit{Slaughter}, the Court in \textit{Trump v.\ Cook} carved out Federal Reserve independence.
\end{enumerate}

\section{Critique}

\textbf{Unitary-executive critique.} Justice Scalia in \textit{Morrison v.\ Olson}, Justice Thomas in \textit{Seila Law}, and scholars such as Steven Calabresi and Saikrishna Prakash argued that Article~II vests all executive power in the President, leaving no home for a headless ``fourth branch,'' and that the FTC exercised real enforcement power anyway. That view prevailed in \textit{Trump v.\ Slaughter}.

\textbf{Historical and functional defense.} Scholars such as Christine Kexel Chabot and Jed Shugerman argue the Founders did not embrace absolute removal power, and that insulating expert, multi-member agencies protects stability and checks executive aggrandizement---as Justice Sotomayor's \textit{Slaughter} dissent warned.

\section{Key Quotations}

\begin{quote}
``The authority of Congress, in creating quasi-legislative or quasi-judicial agencies, to require them to act in discharge of their duties independently of executive control cannot well be doubted, and that authority includes, as an appropriate incident, power to fix the period during which they shall continue in office, and to forbid their removal except for cause in the meantime. For it is quite evident that one who holds his office only during the pleasure of another cannot be depended upon to maintain an attitude of independence against the latter's will.''
\end{quote}
\noindent --- Justice Sutherland, for the Court, 295 U.S.\ at 629.

\begin{quote}
``Such a body cannot in any proper sense be characterized as an arm or an eye of the executive. Its duties are performed without executive leave, and, in the contemplation of the statute, must be free from executive control.''
\end{quote}
\noindent --- Justice Sutherland, for the Court, 295 U.S.\ at 628.

\begin{quote}
``Whether the power of the President to remove an officer shall prevail over the authority of Congress to condition the power by fixing a definite term and precluding a removal except for cause will depend upon the character of the office; the \textit{Myers} decision, affirming the power of the President alone to make the removal, is confined to purely executive officers, and, as to officers of the kind here under consideration, we hold that no removal can be made during the prescribed term for which the officer is appointed except for one or more of the causes named in the applicable statute.''
\end{quote}
\noindent --- Justice Sutherland, for the Court, 295 U.S.\ at 631--32.

\end{document}
